\documentclass[11pt]{amsart}


\usepackage{bm}
\usepackage{amsmath,amsfonts, amsthm, amssymb}
\usepackage[abbrev,backrefs]{amsrefs}
%\usepackage{showlabels}
\usepackage{microtype}
\usepackage{enumerate}
\usepackage{graphicx}
%\usepackage{braket}
\graphicspath{{../../figures/}}
\usepackage[onehalfspacing]{setspace}
\usepackage{nicefrac}
\usepackage{xcolor}
\usepackage{wasysym}

%%%% for widecheck command
% code from mathabx.sty and mathabx.dcl
\DeclareFontFamily{U}{mathx}{\hyphenchar\font45}
\DeclareFontShape{U}{mathx}{m}{n}{
      <5> <6> <7> <8> <9> <10>
      <10.95> <12> <14.4> <17.28> <20.74> <24.88>
      mathx10
      }{}
\DeclareSymbolFont{mathx}{U}{mathx}{m}{n}
\DeclareFontSubstitution{U}{mathx}{m}{n}
\DeclareMathAccent{\widecheck}{0}{mathx}{"71}

\setlength{\oddsidemargin}{0pt}
\setlength{\evensidemargin}{0pt}
\setlength{\textwidth}{6.0in}
\setlength{\topmargin}{0in}
\setlength{\textheight}{8.5in}

\setlength{\parindent}{0in}
\setlength{\parskip}{5px}

\usepackage{scalerel,stackengine}
\def\apeqA{\SavedStyle\sim}
\def\apeq{\setstackgap{L}{\dimexpr.5pt+1.5\LMpt}\ensurestackMath{%
  \ThisStyle{\mathrel{\Centerstack{{\apeqA} {\apeqA} {\apeqA}}}}}}

\def\bra#1{\mathinner{\langle{#1}|}}
\def\ket#1{\mathinner{|{#1}\rangle}}
\def\braket#1{\mathinner{\langle{#1}\rangle}}
\def\Braket#1{\langle{#1}\rangle}

\def\Real{{\mathbf R}}
\def\bbC{{\mathbb C}}
\def\Sphere{{\mathbf S}}
\def\Schwartz{{\mathcal{S}}}
\def\Lebesgue{{\mathcal{L}}}
\def\Hausdorff{{\mathcal{H}}}
\def\Prob{{\mathbf P}}
\def\Var{{\mathbf{Var}\,}}
\def\Expec{{\mathbb E\,}}
\def\xp{{\mathbb E\,}}
\def\RWExpec{{\mathbb E\,}}
\def\eps{{\varepsilon}}
\def\One{{\mathbf{1}}}

\def\PhaseSpace{{\Real^{2d}}}
\def\Evol{{\mathcal{E}}} \def\xop{{\mathbf{x}}}
\def\pop{{\mathbf{p}}}
\DeclareMathOperator{\Rec}{Rec}
\DeclareMathOperator{\Err}{Err}
\def\lsim{{\,\lesssim\,}}
\def\llog{{\,\lessapprox\,}}

% the Hilbert space
\def\Hilbert{{\mathcal{H}}}

\def\bbR{{\mathbb{R}}}
\def\bbZ{{\mathbb{Z}}}
\def\bbN{{\mathbb N}}
\def\bbH{{\mathbb H}}

\newcommand{\Wigner}[1]{{\mathcal{W}_{#1}}}
\newcommand{\Ft}[1]{{\widehat{#1}}}
\newcommand{\IFt}[1]{{\widecheck{#1}}}

\newcommand{\ot}{{\leftarrow}}

\newcommand{\weight}{{\mathrm{weight}}}
\newcommand{\poly}{{\mathrm{poly}}}
\newcommand{\Poly}{{\mathrm{Poly}}}

\newcommand{\diff}{\mathop{}\!\mathrm{d}}

\numberwithin{equation}{section}

\DeclareMathOperator{\Id}{Id}
\DeclareMathOperator{\ad}{ad}
\newcommand{\AD}{{\mathbf{ad}}}
\newcommand{\F}{{\mathcal{F}}}
\newcommand{\eg}{E}
\newcommand{\noset}{\varnothing}
\newcommand{\mbf}[1]{{\mathbf{#1}}}
\newcommand{\mcal}[1]{{\mathcal{#1}}}
\newcommand{\wtild}[1]{{\widetilde{#1}}}
\newcommand{\PathSet}{{\Pi}}
\newcommand{\ddt}{{\frac{d}{dt}}}

\DeclareMathOperator{\Tr}{Tr}

\DeclareMathOperator{\Op}{Op}
\DeclareMathOperator{\Rept}{Re}
\DeclareMathOperator{\Impt}{Im}
\DeclareMathOperator{\vol}{vol}
\DeclareMathOperator{\supp}{supp}
\DeclareMathOperator{\Variance}{Var}
\DeclareMathOperator{\argmax}{argmax}
\DeclareMathOperator{\trace}{tr}
\DeclareMathOperator{\tr}{tr}
\DeclareMathOperator{\Div}{div}

\DeclareMathOperator{\Arm}{arm}
\DeclareMathOperator{\Stick}{stick}

\newtheorem{theorem}{Theorem}
\numberwithin{theorem}{section}
\newtheorem{lemma}[theorem]{Lemma}
\newtheorem{definition}[theorem]{Definition}
\newtheorem{corollary}[theorem]{Corollary}
\newtheorem{proposition}[theorem]{Proposition}
\newtheorem{example}{Example}
\newtheorem{remark}[theorem]{Remark}

\newtheorem{hope}[theorem]{Potentially Provable Lemma}
\newtheorem{assumption}{Useful Assumption}
\hfuzz=20pt

%%%%%%%%

\newcommand{\calD}{\mathcal{D}}
\newcommand{\calP}{\mathcal{P}}
\newcommand{\calS}{\mathcal{S}}
\newcommand{\calY}{\mathcal{Y}}


\newcommand{\nf}{\nicefrac}

\newcommand{\AB}[1]{\textcolor{red}{[AB: #1]}}

\newcommand{\les}{\lesssim}


\newcommand{\bs}{\boldsymbol}

\title{Preliminary hierarchy computations}

\begin{document}
\maketitle

\section{The basics}
I am regurgitating below what I understand from Felix Parraud's notes.

We are going to work on a product space
$\mcal{M}_N(\mcal{A}) = \mcal{C}^{N\times N} \otimes \mcal{A}$ (this is a space
of operators really), so elements of $\mcal{M}_N$ are $N\times N$ matrices where
each operator is an element of the $C^\star$ algebra $\mcal{A}$.  I am thinking here
that $N$ is the number of sites in our periodic lattice $\bbZ^2/n\bbZ^2$ (so $N=n^2$).
The trace on the $\mcal{A}$ is $\tau$, and is normalized so that $\tau(1)=1$,
where $1$ is the identity element.   You can imagine $\mcal{A}$ as being $\mcal{C}^{L\times L}$ where $L\gg N$ is a huge number, and the semicircular elements of
$\mcal{A}$ are GOEs.

Precisely, a free semicircular family $\{s_x\} \subset \mcal{A}$ is a set of elements
which satisfy $s_x^* = s_x$ and for any sequence $(x_1,\cdots,x_k)$ without immediate
repeats ($x_j \not= x_{j+1}$)
\[
\tau(\prod (s_{x_j} - \tau(s_{x_j}))) = 0,
\]
and also
\[
\tau(s_x^k) = \begin{cases} c_{k/2}, &\text{if }k\in2\bbZ, \\ 0, &\text{ else.}\end{cases},
\]
where $c_n$ is the $n$-th Catalan number.

We need the non-commutative derivative $\partial_x$, which is a map
$\partial_x: \mcal{A}\to \mcal{A}\otimes \mcal{A}$.  It is defined so that
\begin{align*}
\partial_x(AB) &= \partial_x A (1\otimes B) + (A\otimes 1) \partial_x B \\
\partial_x s_y &= \delta_{xy} 1\otimes 1 \\
\partial_x (H-z)^{-1} &= - ((H-z)^{-1} \otimes 1) (\partial_x H) (1\otimes (H-z)^{-1}).
\end{align*}
I should say here that we should actually work on some extension $\mcal{F}$ of $\mcal{A}$, which includes $s_x$, all expressions like $A\otimes s_x$ for Hermitian $A\in\mcal{C}^{N\times N}$, and also all inverses like $(H-z)^{-1}$ when $z\in\bbC\setminus \bbR$.  And then in this case the derivative above is intended to
be a definition.  That is, you can sort of think of $\mcal{A}$ as just a list of
formal expressions, and the derivative is a rule for manipulating formal expressions.

What is not a definition but now a \textit{theorem} is the following result.
For $Q\in\mcal{A}$,
\[
\tau(s_x Q(\vec{s})) = (\tau\otimes \tau)(\partial_x Q(\vec{s})).
\]
The trace $\tau\otimes \tau$ on $\mcal{A}\otimes \mcal{A}$ is defined by linearity
and by $(\tau\otimes \tau)(A\otimes B) = \tau(A)\tau(B)$.

\section{Application to Anderson.}
Let $\Delta$ be the Laplacian matrix on a finite torus with $N$ total sites.
Let $\delta_x$ be the diagonal matrix on $\bbC^{N\times N}$ with a $1$ at site $x$.
Make the matrix $H\in\mcal{M}_N(\mcal{A})$
\[
H = \Delta \otimes 1 + \lambda \sum_x \delta_x \otimes s_x =: \Delta + \lambda W
\]
where $\{s_x\}$ is a free semicircular family.
Let's compute the resolvent.  Set
\[
M_\theta := (\Delta - (z+\lambda^2\theta))^{-1}.
\]
Then
\[
(H-z)^{-1} = M_\theta\otimes 1 - (M_\theta\otimes 1)
(\lambda W + \lambda^2\theta) (H-z)^{-1}
\]
(this is an algebra identity valid also on $\mcal{M}_N(\mcal{A})$.
The non-commutative derivative (extended of course to $\mcal{M}_N(\mcal{A})$.
becomes
\[
\partial_x (H-z)^{-1} = -((H-z)^{-1}\otimes 1) (\delta_x\otimes 1) (1 \otimes (H-z)^{-1})
\]
So we'll need to compute
\begin{align*}
\tau(W (H-z)^{-1})
&= \sum_x \delta_x \tau(s_x (H-z)^{-1}) \\
&= \sum_x \delta_x (\tau\otimes\tau)(\partial_x (H-z)^{-1}) \\
&= -\lambda \sum_x \delta_x
(\tau\otimes\tau)((H-z)^{-1}\otimes 1) (\delta_x\otimes 1\otimes 1) (1\otimes  (H-z)^{-1}) \\
&= -\lambda \sum_x \delta_x
\tau((H-z)^{-1})\delta_x \tau((H-z)^{-1}).
\end{align*}
Setting $R = \tau((H-z)^{-1}) \in \bbC^{N\times N}$, the
above becomes
\[
\tau(W(H-z)^{-1}) = -\lambda \mcal{D}[R] R,
\]
so we obtain
\[
R = M_\theta + \lambda^2 M_\theta (\mcal{D}[R]-\theta) R.
\]
Setting $\theta = \mcal{D}[R]$ we get that $R=M_\theta$ with the usual value of $\theta$.

Set now
\[
H_t-z_t := \Delta + \lambda \sqrt{t} W - \eg - (1-t)\lambda^2\theta,
\]
and define $\mcal{R}_t := (H_t-z_t)^{-1}$.
Note that with this definition, $\tau(\mcal{R}_t) = M_\theta$ is constant in time
with the same value of $\theta$.

We are interested in $k$-loops,
\[
K_t(x_1,\cdots,x_k) := \tr \tau(\prod_{j=1}^k \delta_{x_j} \mcal{R}_t(\sigma_j)),
\]
where $\mcal{R}_t(-) := \mcal{R}_t^*$ and $\mcal{R}_t(+) := \mcal{R}_t$.  For
now let's drop that notation because it shouldn't matter for the algebra.

So I want to compute $\partial_t K_t(\vec{x})$.  Note that
\[
\partial_t \mcal{R}_t = - \mcal{R}_t (\frac{\lambda}{2\sqrt{t}}W + \lambda^2\theta)
\mcal{R}_t.
\]
Later we'll also need
\[
\partial_z \mcal{R}_t = -\lambda \sqrt{t} (\mcal{R}_t \otimes 1) \delta_z (1\otimes \mcal{R}_t).
\]
So:
\begin{align*}
\partial_t K_t(\vec{x})
= -\sum_{j_0=1}^k
\tr \tau( (\prod_{j<j_0} \delta_{x_j}\mcal{R}_t)
\delta_{x_{j_0}} \mcal{R}_t (\frac{\lambda}{2\sqrt{t}} W + \lambda^2\theta)\mcal{R}_t
(\prod_{j>j_0} \delta_{x_j} \mcal{R}_t)).
\end{align*}
Let's compute the term with $j_0=1$ (by the cyclic property of the trace this is the
only case to worry about anyway):
\begin{align*}
-\tr \tau (\delta_{x_1} \mcal{R}_t (\frac{\lambda}{2\sqrt{t}}W + \lambda^2 \theta)
\mcal{R}_t (\prod_{j>1} \delta_{x_j}\mcal{R}_t)) &=
- \sum_z \tr\tau(\delta_{x_1} \mcal{R}_t (\frac{\lambda}{2\sqrt{t}}\delta_z\otimes s_z
+ \lambda^2 \theta) \mcal{R}_t (\prod_{j>1} \delta_{x_j} \mcal{R}_t))
\end{align*}
We are going to integrate by parts in $s_z$, using the product rule
(you can use the cyclic property of the trace again to put the $s_z$ on the left if you
prefer, to make the integration by parts and the product rule more familiar).
Let's first worry about the derivative that hits the $\mcal{R}_t$ all the way on the
left:
\begin{align*}
\frac{\lambda}{2\sqrt{t}}
\sum_z
-\tr (\tau\otimes\tau) (\delta_{x_1} &(\partial_z\mcal{R}_t) \delta_z
((\mcal{R}_t (\prod_{j>1} \delta_{x_j}\mcal{R}_t))\otimes 1))  \\
&= \frac{\lambda^2}{2} \sum_z \tr(\tau\otimes \tau)
(\delta_{x_1}(\mcal{R}_t\otimes 1) \delta_z (1\otimes \mcal{R}_t) \delta_z (\mcal{R}_t (\prod_{j>1} \delta_{x_j}\mcal{R}_t))\otimes 1) \\
&= \frac{\lambda^2\theta}{2} \sum_z \tr \tau(\delta_{x_1}\mcal{R}_t \delta_z \mcal{R}_t
\prod_{j>1} \delta_{x_j} \mcal{R}_t).
\end{align*}
The same exact term appears when the derivative hits the $\mcal{R}_t$ immediately to
the right of $\theta$, and this cancels the $\theta$.  Therefore, using
the cylic property of the trace to rearrange the remaining
\begin{align*}
-\tr \tau (\delta_{x_1} \mcal{R}_t &(\frac{\lambda}{2\sqrt{t}}W + \lambda^2 \theta)
\mcal{R}_t (\prod_{j>1} \delta_{x_j}\mcal{R}_t))\\
&=-\frac{\lambda}{2\sqrt{t}}\sum_z
\tr(\tau\otimes \tau)
((\mcal{R}_t\otimes 1) \partial_z(\prod_{j>1} \delta_{x_j}\mcal{R}_t) \delta_{x_1}(1\otimes \mcal{R}_t)) \\
&= \frac{\lambda^2}{2}
\sum_z \sum_{j'>1} \tr(\tau\otimes \tau)
((\mcal{R}_t\otimes 1) (\prod_{1<j<j'} \delta_{x_j}\mcal{R}_t \otimes 1)
\delta_{x_{j'}}
(\mcal{R}_t\otimes 1) \delta_z (1\otimes \mcal{R}_t) (\prod_{j>j'} \delta_{x_j}\mcal{R}_t) ( 1\otimes \mcal{R}_t)) \\
&= \frac{\lambda^2}{2} \sum_z \sum_{j'>1}K_t(x_1,z,x_{j'+1},\dots,x_k)
K_t(z,x_2,\cdots,x_{j'}).
\end{align*}
This accounts for the cuts \textit{from} the $x_1$-$x_2$ edge, and summing over all
edges gives us all the cuts (and the factor of $\frac12$ accounts for the double-counting).


\end{document}
